\subsection*{一、研究動機}
光散射（light scattering）法不需要前處理，元件也便宜，常用來分析水中微粒的大小，
例如微塑膠（microplastics）。多角度光散射利用不同角度的散射強度辨識粒徑，但多半在
稀薄的樣品中進行：此時光大多只被散射一次，角度分布就是單一粒子的散射圖形。樣品變濃
後，光會在樣品中被散射很多次（多重散射），方向逐漸被打亂，角度分布也跟著變平。我們
想知道：從單次散射過渡到多重散射時，散射光的角度分布會怎麼改變？兩個角度的強度比是
否仍與粒徑有關？蒙地卡羅模擬能否重現實驗的結果？

\subsection*{二、研究目的}
\begin{itemize}
    \item[(一)] 建立低成本的多角度光散射量測系統與數據處理流程。
    \item[(二)] 探討光學厚度 $\tau$ 對兩種粒徑（$0.49 \mu m$、$2.16 \mu m$）聚苯乙烯懸浮液兩角強度比 $C$ 的影響。
    \item[(三)] 比較實驗與蒙地卡羅模擬的結果。
\end{itemize}
\subsection*{三、文獻回顧}
已有研究指出，微塑膠的角度散射具有粒徑特徵（Genc et al., 2023），也有研究以多角度光散射辨識聚苯乙烯微球粒徑（Lotfi Choobbari et al., 2022）；然而，這些研究主要在稀薄、以單次散射為主的條件下進行。另一方面，蒙地卡羅模擬已用於描述聚苯乙烯懸浮液中的多重散射（Berrocal et al., 2007）。對於單次散射過渡至多重散射的區間，兩個角度的強度比如何隨光學厚度改變，尚未見直接相關研究。因此，本研究以兩角強度比作為指標，探討其在不同粒徑與光學厚度下的變化，並以蒙地卡羅模擬與實驗互相比較。
